% /root/monodromy, 2026-11-19. Complete root, reduction-descent and obstructions-crystals audits passed.
\subsection{Uniform tangent triples for odd graph kernels}
\label{subsec:ss-uniform-tangent-triples}
\status{Proved and independently checked. The parent, reduction/descent and
obstructions/crystals agents independently audited the complete geometric
proof and both dependency files. The actual rank-five witness certificate
was independently read and rerun. This proves uniform tangent-triple
witnesses for all odd graph kernels; its full-group consequence is stated
separately.}

Continue with $E:y^2+y=x^3$, $X=E^2$, $L=\mathcal O(3O)\boxtimes
\mathcal O(3O)$ and a finite odd subgroup $K\subset X(k)$ whose
two coordinate projections are injective. Given a triple of distinct
elements, translate and reorder it as $\{0,g,\eta\}$.
We may and do arrange
\[
 \eta\ne-g,\qquad\hbox{or}\qquad 3g=0\ \hbox{and}\ \eta=-g.
\]
Indeed if the original triple is $\{0,g,-g\}$ and $3g\ne0$, use
the translated triple $\{0,2g,g\}$. Take the contact space
$\Lambda$ defined by $g$, and its third-value hyperplane
$H_\eta=(l_\eta=0)\simeq\mathbf P^3$, where $l_j(s)=s(j)$.
The coordinates $F,T,U,v,h$, functions $X_1,Y_1,Z_1,\ell,P_1$,
and map $\Phi$ are those of
Lemma~\ref{lem:ss-contact-cubic-evaluation}.

\begin{lemma}[The projected value map and its three-value ranks]
\label{lem:ss-third-value-projection}
If $\eta\ne-g$, the base locus of $H_\eta$ consists of
$P=(O,O),Q=g,\eta$, and the possible additional point
$(g_1,-2g_2)$ already in the contact base locus. The additional
point has no $K$-difference with the other three.
For every $0\ne j\in K$, the two values at $R,R+j$ on $H_\eta$
have rank one only at finitely many $R$, and generically have rank
two. For every three distinct translates their three value
functionals are generically independent.

If $3g=0$ and $\eta=-g$, the base locus consists of $P,Q,-g$.
The pair-rank assertion still holds for $j\notin\{0,g,-g\}$.
For $j=\pm g$, any additional positive-dimensional rank-one
stratum contributes only sections with $F=0$; it can be discarded
on $F\ne0$. The generic three-value independence holds in this
case as well.
\end{lemma}
\begin{proof}
The new value map is projection of $\Phi$ from the point
$\Phi(\eta)$. Away from the exceptional fibres in
Lemma~\ref{lem:ss-contact-cubic-evaluation}, $\Phi$ is injective.
For $\eta\ne\pm g$, its only possible exceptional horizontal
fibre is $\eta=-2g$, since $\ell(\eta)=0$ exactly there.
On that fibre the map is the degree-two pencil
\[
                         \rho(x,y)=\frac{y+b}{x+a}.
\]
The fibre through $-2g_1$ has other point $g_1$, which is already
the base point $(g_1,-2g_2)$. Hence no further base point appears.
If $\eta=-g$ and $3g=0$, the value is $F+T$; its fibre on the
vertical exceptional line is given by $\ell=1$, whose support
is just $-g_2$ because this is a flex. These observations give
the stated base loci. The assertion about $K$-differences follows
from injectivity of the first projection, as before.

We make the pair-rank assertion precise. Where the first coordinates
of $R,R+j$ avoid $O,\pm g_1$, divide the five value coordinates
by $X_1$. They become
\begin{equation}\label{eq:ss-projected-value-coordinates}
              (b_0,\rho b_0,b_2/X_1,b_1,b_2),
\end{equation}
where $(b_0,b_1,b_2)$ is the homogeneous value vector of the
second elliptic point in the basis $(1,z,\ell)$.
Three value rows at $R,R+j,\eta$ can be dependent only if their
projections to $(b_0,b_1,b_2)$ are dependent. If
$\eta\ne-g$, its projection is nonzero. The three second-factor
points must therefore be collinear, unless two coincide. The
noncoincident case gives
\[
                            2R_2+j_2+\eta_2=0,
\]
so there are only finitely many possibilities for $R_2$.
Coincidences also give only finitely many possibilities.

Fix one of these second-factor values. In a dependence relation,
the second coordinate in
\eqref{eq:ss-projected-value-coordinates} would give a
nontrivial relation
\[
                  A\rho(R_1+j_1)+B\rho(R_1)+C=0.
\]
There is no such identity in $R_1$. The pole divisor of $\rho$
is exactly $(O)+(-g_1)$, with two simple poles. If both $A,B$
are nonzero, an identity would make this divisor invariant under
translation by $j_1$. Its stabilizer has order dividing two,
whereas $j_1$ has nontrivial odd order. If one coefficient is zero,
the nonconstant remaining function gives the contradiction.
The case in which a second-factor value is $O$ uses the same
homogeneous formula: the other moving value has $b_0\ne0$, and
its coefficient cannot be zero unless the fixed second point
were $O$, which it is not. Thus the relation remains nontrivial.
We conclude that the rank-one points on these finitely many
horizontal fibres are finite.

The omitted first-factor exceptional fibres cause no new curve
when $\eta\ne-g$. Their projections to $(b_0,b_1,b_2)$ are zero.
If only one moving point lies on such a fibre, dependence forces
the other second point to equal $\eta_2$, giving finitely many
points. If both do, dependence reduces to equality of their
original $\Phi$-values, already proved to occur at only finitely
many points. For $\eta=-g$, the fixed projection is zero.
On the nonexceptional open dependence would require the two
distinct second-factor points to have proportional value vectors,
which is impossible. Two exceptional vertical fibres can give
a curve only when $j=\pm g$ (since $3g=0$). Their values lie on
the line spanned by $F,T$; on $H_{-g}$ one has $T=F$, so a
generic point of every such incidence requires $F=0$.
All remaining exceptional points are finite. This proves the
pair assertions, including their stated restriction in the cyclic case.

We finally prove generic independence for three values on
$H_\eta$. If they were dependent, the fixed value
$\Phi(\eta)$ would lie in the span of the three original values
at $R,R+j,R+k$ for a geometric generic $R$.
The second-factor value vectors of these three points are generically
a basis: dependence would impose $3R_2+j_2+k_2=0$.
If $\eta=-g$, its projection is zero, forcing all three coefficients
to vanish, which is impossible.
Otherwise normalize the first coordinate of each second-factor
value to one. The coefficients $\alpha_i\in k(R_2)$ expressing
the fixed second-factor point have sum one. Their other value
equation would force
\begin{equation}\label{eq:ss-three-pole-relation}
       \alpha_0\rho(R_1)+\alpha_1\rho(R_1+j_1)
             +\alpha_2\rho(R_1+k_1)=\rho(\eta_1).
\end{equation}
The poles of these three functions are the three pairs
$\{-j_i,-g_1-j_i\}$, with $j_i=0,j_1,k_1$.
Every nonzero coefficient must have all its poles cancelled.
Viewing the pairs as edges in orbits of translation by $g_1$,
a set of at most three distinct edges can have no endpoint of
degree one only if it is an entire three-cycle. Thus the only
possible case is $3g=0$ and $\{0,j,k\}=\{0,g,-g\}$.
The two residues of $\rho$ against an invariant differential are
equal and nonzero (by the residue theorem in characteristic two).
Cancellation in this cycle therefore forces all $\alpha_i$ equal,
and their sum one forces each to be one. But the sum of the three
second-factor $x$-coordinates is not constant: it has a pole at
each of the three distinct translated origins, with no cancellation.
It cannot equal the fixed $x(\eta_2)$ for generic $R_2$.
This contradicts \eqref{eq:ss-three-pole-relation}, and
completes the proof.
\end{proof}

\begin{lemma}[Two linear independence tests]
\label{lem:ss-third-value-linear-tests}
Let $i,\eta\in K\setminus\{0,g\}$ be distinct, with
$i\ne-g$. On $\ker(l_i,l_\eta)\subset\Lambda$, neither $h$ nor
$\partial_2s(i)$ is identically zero.
If additionally $i=2g$, $3g\ne0$ and $\eta\ne\pm g,2g$,
then $v$ is not identically zero there either.
\end{lemma}
\begin{proof}
For a point $j\ne\pm g$ write
\[
 X_j=x(j_1)+a,\quad Z_j=x(j_2)+c,\quad
 p_j=\rho(j_1)+a^2,\quad r_j=\ell(j)/X_j.
\]
The normalized value row is
\[
 l_j/X_j=p_jF+r_jT+Z_jU+v+\ell(j)h,
\]
whereas
\[
 \partial_2s(j)/X_j=(Z_j^2/X_j)T+U+Z_j^2h.
\]
The rows of $l_i,l_\eta$ modulo $F,h$ are independent, by the
same argument with $(\ell(j),X_j,X_jZ_j)$ used in
Lemma~\ref{lem:ss-contact-no-triples}. This proves the assertion
for $h$. If $\eta=-g$, then $l_\eta=F+T$, and independence of
the derivative follows at once from its nonzero $U$ coefficient
and the nonzero $v$ coefficient of $l_i$.

Assume $\eta\ne\pm g$. A relation expressing the derivative in
the span of the two value rows must cancel their $v$ coefficients.
Its $F$ coefficient then requires $p_i=p_\eta$.
The degree-two pencil $\rho$ has fibre divisors summing to $-g_1$;
as $i\ne\eta$, this would imply $i+\eta=-g$ by injectivity of
the first projection. Hence the three distinct second-factor points
$i_2,\eta_2,g_2$ lie on one affine line, of slope $\beta$.
Put $L_0=\beta+c^2$. Then
\[
 \ell(i)=L_0Z_i,\qquad \ell(\eta)=L_0Z_\eta,
 \qquad Z_i+Z_\eta\ne0.
\]
Comparing the $U$ and $h$ coefficients in the alleged relation
gives $L_0=Z_i^2$, or $\beta=x(i_2)^2$.
Thus that line would be tangent at $i_2$, despite already containing
the three distinct points $i_2,\eta_2,g_2$. This is impossible.
The derivative assertion follows.

For the final assertion, $i=2g$ has $\ell(i)\ne0$: the zeros of
$\ell$ are $g_2,-2g_2$, neither equal to $2g_2$ for nonzero odd
$g$ with $3g\ne0$. If $v$ lay in the span of $l_i,l_\eta$, the
four coordinates other than $v$ in the normalized rows would be
proportional. Their $T$ and $h$ entries force $X_i=X_\eta$.
Hence $\eta=-i$ (the other possibility $\eta=i$ is excluded).
Their nonzero $U$ entries then force proportionality scalar one,
but $\ell(-i)=\ell(i)+1$, a contradiction. This proves the last
claim.
\end{proof}

\begin{lemma}[No unwanted collisions on the third-value slice]
\label{lem:ss-third-value-other-collisions}
Suppose first that $\eta\ne-g$. A geometric generic section of
$H_\eta$ has exactly its prescribed three-point orbit and no
other triple identifications under $K$. All pair identifications
are transverse except the pair $P,Q$. Its contact length is two
if $3g\ne0$ and three if $3g=0$.

If $3g=0$ and $\eta=-g$, a generic section of $H_\eta$ has no
nontransverse pair of difference outside $\langle g\rangle$,
and no triple with at least one pairwise difference outside that subgroup.
No assertion about its prescribed three tangent directions is
made in this second case.
\end{lemma}
\begin{proof}
For each $j\ne0$ use the proper two-value incidence over $H_\eta$.
Its generic coefficient fibre over the moving point $R$ is a
projective line. The projected-value lemma shows that there is
one integral three-dimensional main component, and that all other
components dominating $H_\eta$ are the constant base-point pairs.
In the cyclic case the additional curve strata lie over $F=0$,
which we remove. The intersection family is generically finite,
and its total length is $L^2=18$.

First assume $\eta\ne-g$ and $j\ne\pm g$.
If $j\ne\eta$, set $R=P$ and impose $l_j=0$ as well as
$l_\eta=0$. By the linear tests choose $h\partial_2s(j)\ne0$.
The two value equations have invertible Jacobian in $R$, so this
is an etale point of the main component. If $j=\eta$, instead
take $R=-\eta$, impose $l_{-\eta}=0$, and use the same test
with the constant second point $P$. In both cases the moving
component is dominant and generically etale.

For $j=g$ with $3g\ne0$ and $\eta\ne2g$, take $R=Q$.
The two value equations are the prescribed zero at $Q$ and the
additional condition $l_{2g}=0$; the last linear test permits
$v\partial_2s(2g)\ne0$. This again gives an etale point of the
main component. If instead $\eta=2g$, the constant pairs have
lengths two at $P$ and one at $Q$: the third branch has nonzero
second-factor derivative, as its derivative functional has nonzero $U$
coefficient whereas $l_\eta$ has nonzero $v$ coefficient.
The moving component therefore has degree $18-3=15$, which is
odd and hence separable in characteristic two.
If $3g=0$ (still $\eta\ne-g$), the sole constant $g$-pair is
$P,Q$ with length three. Its coefficient $T+vF/h$ is not identically
zero on $H_\eta$. Indeed its homogeneous numerator $hT+vF$ is an
irreducible quadratic in the independent variables $F,T,v,h$,
so no hyperplane divides it. The moving degree is again $15$.
The contact length is two in the nonflex case by
Lemma~\ref{lem:ss-uniform-local-contact}. The assertion for $-g$
is its translate. All other constant pairs involve the third point
and are transverse on a dense open, by the nonzero second-factor derivative
there and the nonzero first derivatives $h,v$ at $P,Q$.
Thus after shrinking the coefficient base all moving intersections
are etale and the only nontransverse prescribed pair is $P,Q$.

In the cyclic case $H_{-g}$ has $T=F$.
For every $j\notin\{0,g,-g\}$, take $R=P$ and impose $l_j=0$.
The linear tests give $h\partial_2s(j)\ne0$, proving that the
unique moving component is generically etale. There is no constant
base pair of such a difference.

For triples, the projected-value lemma gives generic rank three;
the coefficient fibre is then a point. The rank-two stratum has
base dimension at most one and coefficient fibres of dimension
one. Rank-one strata are finite with coefficient fibres of dimension
two; in the cyclic exceptional case their curve parts require
$F=0$ and have been removed. Thus every nonconstant triple
incidence has dimension at most two, and cannot dominate
$H_\eta\simeq\mathbf P^3$. The only rank-zero triples are the
prescribed three base points. This proves precisely the assertions
about triples in both cases. There are finitely many differences
and triples, so all stated open restrictions can be imposed at once.
\end{proof}

\begin{proposition}[A tangent-triple bridge for every kernel triple]
\label{prop:ss-uniform-tangent-triple}
Let $K\subset E^2(k)$ be a finite odd subgroup with injective
coordinate projections. For every triple of distinct elements of
$K$, there exists a smooth sharp source $C\in|L|$ whose image
in $X/K$ has one singularity with those three preimages: two
branches have simple contact and the third is transverse.
Every other singularity is an ordinary node.
Thus the already proved local tangent-triple Cartier criterion
applies to every triple-collision incidence.
\end{proposition}
\begin{proof}
Use the ordering $\{0,g,\eta\}$ fixed above.
Lemma~\ref{lem:DE-A4-uniform-triple-local-type} constructs the
irreducible main rank-five triple-contact incidence, a
$\mathbf P^3$-bundle over an open of the moving point and tangent
base. It contains our vertical slice $H_\eta$.
That lemma and Corollary~\ref{cor:DE-A4-third-value-smooth} give
a nonempty open of globally smooth sharp sources with two simply
tangent branches and a transverse third branch.

If $\eta\ne-g$, the preceding lemma supplies another nonempty
open, excluding every unwanted collision. We justify openness
at its prescribed points. The third branch is transverse already
on the vertical slice, so the three-value incidence is reduced
at the prescribed triple: two of its value derivatives are
independent. Thus an extra triple cannot specialize into it.
All constant transverse pairs are likewise etale. The only remaining
cluster is the prescribed tangent pair, of length two or three.
In the full contact incidence it always has length at least two;
a second tangency in that cluster would require length at least
four. The finite flat intersection argument of
Proposition~\ref{prop:ss-uniform-graph-mixed-contact} therefore
gives the required open even when $3g=0$.
Intersect the two nonempty opens in the irreducible main incidence.

For $3g=0$ and $\eta=-g$, apply the preceding lemma only to
differences outside $\langle g\rangle$ and triple patterns with at least
one pairwise difference outside $\langle g\rangle$.
They have no prescribed base pairs or triples, so their avoidance
is open directly by the proper incidence maps.
For the collisions inside $\langle g\rangle$, use the already
certified collinear tangent-triple witness of
Lemma~\ref{lem:ss-global-tangent-triples}.
Product automorphisms preserving $O$ carry its nonzero three-torsion
generator to $g$ in both coordinates. The resulting witness has
exactly one transverse-third tangent triple and no other
$\langle g\rangle$-collision. This supplies a nonempty open of
the same main incidence. Its four contact rows and its third value
have rank five, as checked at the actual witness in the retained
program
\path{agents/monodromy/code/verify_ss_tangent_incidence_rank.py}.
All source, contact and internal-collision conditions are open at
this witness; the transverse third makes the prescribed triple
reduced. Irreducibility lets this open intersect all of the
outside-collision opens obtained from the vertical slice.

In either case translated intersections are finite, so the smooth
source maps birationally to its image. The collision conditions
give exactly the asserted singularities. Sharpness is
$\det S\,Q_H(S)\ne0$, equivalently the only Cartier-killed source
forms are the two invariant ambient forms. The local conclusion
is Lemma~\ref{lem:ss-tangent-triple-bridge}. No classification of
the resulting full monodromy group is included in this proposition.
\end{proof}
